% SYNTHETIC TEST UNIT (integration tools) -- deliberate errors are marked "DELIBERATE"
\chapter{总结（合成测试）}\label{ch:summary}

\section{主表}\label{sec:summary-master}

% DELIBERATE C-P11 E (row II.2 prints the wrong tag; row III.4 names an undefined label); C-P14 W (row III.2)
\begin{table}[htbp]\centering\footnotesize
\caption{全书主要结论总表（合成测试）}\label{tab:summary-master}
\begin{tabular}{@{}lllllll@{}}\toprule
\# & 结论 & 标签 & 非微扰地位 & 规范标签 & 瓶颈假设 & ★ \\\midrule
II.1 & 矩的因子化 & \tagR\ $R_{\mathrm{OPE}}$ & W & \texttt{thm:dis-moments} & \cref{thm:dis-ope} & ★ \\
II.2 & 欧氏 OPE & \tagR\ $R_{\mathrm{PT}}$ & 部分子 & \texttt{thm:dis-ope} & — & \\
III.1 & DY 因子化 & \tagA & Br & \cref{thm:dy-factorization} & \cref{asm:appN-collins} & ★ \\
III.2 & 模型命题 & \tagP & 部分子 & \texttt{prop:dy-model} & \cref{asm:appN-unused} & \\
III.3 & 强子层面命题 & \tagA & Br & \texttt{prop:dy-hadron} & \cref{asm:appN-collins} & \\
III.4 & 不存在的定理 & \tagA & — & \texttt{thm:dis-nonexistent} & — & \\
III.5 & 条件减除定理 & \tagR\ $R_{\mathrm{cond}}$ & 部分子 & \texttt{thm:appN-subtraction} & \cref{asm:appN-collins} & ★ \\
\bottomrule\end{tabular}\end{table}
